\documentclass[aspectratio=169]{beamer}

% Tema UFS --- Universidade Federal de Sergipe
\usetheme{UFS}

% Pacotes base
\usepackage[utf8]{inputenc}
\usepackage[portuguese]{babel}
\usepackage{amsmath,amssymb}
\usepackage{graphicx}
\usepackage{booktabs}
\usepackage{xcolor}

% TikZ e bibliotecas
\usepackage{tikz}
\usetikzlibrary{shapes.geometric, shapes.multipart, arrows.meta, positioning,
  matrix, fit, backgrounds, decorations.pathreplacing, patterns, calc}

% Listagens --- paleta VS Code Light+
\usepackage{listings}
\definecolor{bgcolor}{HTML}{FFFFFF}
\definecolor{linecolor}{HTML}{DDDDDD}
\definecolor{numcolor}{HTML}{999999}
\definecolor{kwcolor}{HTML}{0000FF}
\definecolor{strcolor}{HTML}{A31515}
\definecolor{cmtcolor}{HTML}{008000}

\lstdefinestyle{ufsStyle}{
  backgroundcolor=\color{bgcolor},
  basicstyle=\ttfamily\footnotesize\color{black},
  keywordstyle=\color{kwcolor}\bfseries,
  stringstyle=\color{strcolor},
  commentstyle=\color{cmtcolor}\itshape,
  numberstyle=\tiny\color{numcolor},
  numbers=left, numbersep=12pt,
  xleftmargin=20pt, framexleftmargin=20pt,
  breaklines=true, tabsize=4,
  keepspaces=true, showspaces=false,
  showstringspaces=false, showtabs=false,
  frame=single, rulecolor=\color{linecolor},
  framesep=4pt, captionpos=b, upquote=true,
  columns=fullflexible,
}
\lstset{style=ufsStyle, language=C}

\lstdefinestyle{tinyC}{
  style=ufsStyle, language=C, numbers=none,
  basicstyle=\ttfamily\scriptsize\color{black},
}
\lstdefinestyle{codC}{
  style=tinyC, basicstyle=\ttfamily\fontsize{7.5}{8.6}\selectfont\color{black},
}
\lstdefinestyle{shellStyle}{
  style=ufsStyle, language=bash, numbers=none,
  basicstyle=\ttfamily\scriptsize\color{black},
}

% --- Estilos TikZ reutilizados ---
\tikzset{
  no/.style={rectangle split, rectangle split horizontal, rectangle split parts=2,
    draw=UFSBlue, fill=UFSBlue!8, minimum height=0.55cm,
    font=\scriptsize\ttfamily, inner sep=3pt,
    rectangle split part fill={UFSBlue!8, white}},
  nonovo/.style={no, draw=UFSRed, rectangle split part fill={UFSRed!8, white}},
  nolivre/.style={no, draw=UFSGray, dashed, rectangle split part fill={black!5, black!5}},
  var/.style={rectangle, draw=UFSGray, fill=black!5, minimum width=0.7cm,
    minimum height=0.55cm, font=\scriptsize\ttfamily},
  cel/.style={rectangle, draw=UFSBlue, fill=UFSBlue!8, minimum width=0.62cm,
    minimum height=0.5cm, font=\scriptsize\ttfamily, inner sep=1pt},
  celnova/.style={cel, draw=UFSRed, fill=UFSRed!12},
  nulo/.style={font=\scriptsize\ttfamily, text=UFSGray},
  rot/.style={font=\scriptsize\ttfamily, text=UFSGray},
  leg/.style={font=\scriptsize, text=UFSGray},
  seta/.style={-{Stealth[length=2mm]}, thick, UFSBlue},
  setar/.style={-{Stealth[length=2mm]}, thick, UFSRed},
  setag/.style={-{Stealth[length=2mm]}, thick, UFSGray},
}

% Centro da parte "prox" de um no (ponto de partida da seta)
\newcommand{\ptr}[1]{($(#1.one split north)!0.5!(#1.south east)$)}

% Referencia da fonte no rodape do frame
\newcommand{\fonte}[1]{\vfill\hfill{\scriptsize\textcolor{UFSGray}{Fonte: #1}}}

% --- Metadados ---
\title{Listas encadeadas}
\subtitle{Uma introdução: o nó, o encadeamento e o caminho até o \texttt{NULL}}
\author{Prof.\ André Yoshiaki Kashiwabara}
\institute{DCOMP -- Universidade Federal de Sergipe\\
\texttt{andreyoshiaki@dcomp.ufs.br}}
\date{}

\begin{document}

\begin{frame}[plain]
  \titlepage
\end{frame}

\begin{frame}{O que você vai aprender hoje}
  \begin{center}
  \begin{tikzpicture}[
    item/.style={rectangle, draw=UFSBlue, fill=UFSBlue!10, rounded corners,
      minimum width=10.4cm, minimum height=0.55cm, font=\small}
  ]
    \node[item] (t1) at (0, 0)    {Por que não basta um vetor};
    \node[item] (t2) at (0,-0.72) {O nó: um registro que aponta para outro};
    \node[item] (t3) at (0,-1.44) {Montar e percorrer: do início até o \texttt{NULL}};
    \node[item] (t4) at (0,-2.16) {Inserir no início};
    \node[item] (t5) at (0,-2.88) {Liberar a lista sem perder o caminho};
    \node[item] (t6) at (0,-3.60) {Vetor ou lista?};
    \node[item] (t7) at (0,-4.32) {Exercícios};
    \draw[seta] (t1) -- (t2);
    \draw[seta] (t2) -- (t3);
    \draw[seta] (t3) -- (t4);
    \draw[seta] (t4) -- (t5);
    \draw[seta] (t5) -- (t6);
    \draw[seta] (t6) -- (t7);
  \end{tikzpicture}
  \end{center}
\end{frame}

% ==================================================================
\section{A pergunta de hoje}
% ==================================================================

\begin{frame}{O problema: guardar uma coleção que muda o tempo todo}
  \begin{columns}[T]
    \begin{column}{0.5\textwidth}
      \textbf{Até aqui}\\[0.3em]
      \small
      A \texttt{agenda.c} guardava os contatos num vetor no heap, que
      crescia com \texttt{realloc}. Funciona --- enquanto os contatos
      entram \emph{no fim}.

      \medskip
      E se cada contato novo precisar entrar \textcolor{UFSRed}{no
      começo}? Ou no meio, em ordem alfabética?
    \end{column}
    \begin{column}{0.5\textwidth}
      \textbf{O que vamos descobrir}\\[0.3em]
      \small
      Um vetor exige que os elementos fiquem \textbf{lado a lado} na
      memória. Abrir mão disso --- cada elemento num bloco próprio,
      apontando para o próximo --- é a ideia da \textbf{lista
      encadeada}, a primeira estrutura de dados dinâmica do curso.
    \end{column}
  \end{columns}
  \begin{block}{A pergunta de hoje}
    Como guardar uma sequência sem que os elementos precisem estar
    vizinhos na memória --- e o que isso custa?
  \end{block}
\end{frame}

\begin{frame}[fragile]{O programa de hoje: \texttt{lista.c} (1/3)}
  \begin{columns}[T]
    \begin{column}{0.46\textwidth}
\begin{lstlisting}[style=codC]
/* lista.c --- uma lista encadeada */
#include <stdio.h>
#include <stdlib.h>

typedef struct no {
    int valor;
    struct no *prox;  /* ou NULL */
} No;
\end{lstlisting}
      \small Cada valor da lista mora num \texttt{No}, alocado com
      \texttt{malloc}.
    \end{column}
    \begin{column}{0.5\textwidth}
\begin{lstlisting}[style=codC]
/* devolve o novo inicio da lista */
No *insere_inicio(No *ini, int valor)
{
    No *novo = malloc(sizeof *novo);
    if (novo == NULL) {
        perror("malloc");
        exit(1);
    }
    novo->valor = valor;
    novo->prox = ini;
    return novo;
}
\end{lstlisting}
    \end{column}
  \end{columns}
\end{frame}

\begin{frame}[fragile]{\texttt{lista.c} (2/3): percorrer e liberar}
  \begin{columns}[T]
    \begin{column}{0.48\textwidth}
\begin{lstlisting}[style=codC]
void imprime(const No *ini)
{
    for (const No *p = ini; p != NULL;
         p = p->prox)
        printf("%d -> ", p->valor);
    printf("NULL\n");
}

int conta(const No *ini)
{
    if (ini == NULL)
        return 0;
    return 1 + conta(ini->prox);
}
\end{lstlisting}
    \end{column}
    \begin{column}{0.48\textwidth}
\begin{lstlisting}[style=codC]
void libera(No *ini)
{
    while (ini != NULL) {
        No *prox = ini->prox; /* antes! */
        free(ini);
        ini = prox;
    }
}
\end{lstlisting}
    \end{column}
  \end{columns}
\end{frame}

\begin{frame}[fragile]{\texttt{lista.c} (3/3): o \texttt{main}}
  \begin{columns}[c]
    \begin{column}{0.52\textwidth}
\begin{lstlisting}[style=codC]
int main(void)
{
    No *lista = NULL;          /* vazia */
    for (int i = 1; i <= 4; i++)
        lista = insere_inicio(lista, 10 * i);
    imprime(lista);
    printf("%d nos, %zu bytes cada\n",
           conta(lista), sizeof(No));
    libera(lista);
    return 0;
}
\end{lstlisting}
    \end{column}
    \begin{column}{0.44\textwidth}
      \small
      A lista começa vazia: \texttt{lista} vale \texttt{NULL}.

      \medskip
      Cada volta do laço cria um nó com o valor $10 \cdot i$ e guarda em
      \texttt{lista} o ponteiro que \texttt{insere\_inicio} devolve.
    \end{column}
  \end{columns}
\end{frame}

\begin{frame}{O que este programa imprime?}
  \begin{block}{Antes de compilar, escreva sua resposta}
    \begin{enumerate}
      \item Em que ordem os valores 10, 20, 30 e 40 aparecem na primeira
        linha?
      \item Quantos nós \texttt{conta} encontra?
      \item Quantos bytes tem cada \texttt{No}? Um \texttt{int} tem 4 bytes
        e um ponteiro, 8.
      \item Quantas vezes \texttt{malloc} é chamado? E \texttt{free}?
    \end{enumerate}
  \end{block}
\end{frame}

\begin{frame}[fragile]{Compile, execute e compare}
\begin{lstlisting}[style=shellStyle]
$ gcc -Wall -Wextra lista.c -o lista
$ ./lista
40 -> 30 -> 20 -> 10 -> NULL
4 nos, 16 bytes cada
\end{lstlisting}
  \begin{columns}[T]
    \begin{column}{0.48\textwidth}
      \begin{block}{Duas surpresas}
        \small
        \begin{itemize}
          \item a ordem saiu \textbf{invertida};
          \item 16 bytes, e não $4 + 8 = 12$.
        \end{itemize}
      \end{block}
    \end{column}
    \begin{column}{0.48\textwidth}
      \small Quatro \texttt{malloc}, um por nó, e quatro \texttt{free}. Ao
      longo da aula, cada surpresa ganha uma explicação.
    \end{column}
  \end{columns}
\end{frame}

% ==================================================================
\section{Por que não basta um vetor}
% ==================================================================

\begin{frame}{Inserir no começo de um vetor custa quanto?}
  \begin{columns}[c]
    \begin{column}{0.5\textwidth}
      \begin{center}
      \begin{tikzpicture}
        \foreach \x/\v in {0/30,1/20,2/10} {
          \node[cel] (a\x) at (0.66*\x+0.66, 0) {\v};
        }
        \node[cel, fill=white] at (2.64, 0) {};
        \node[celnova] (n) at (0, 0.9) {40};
        \draw[setar] (n) -- (0, 0.25) -- (0.33,0.25);
        \foreach \x in {0,1,2} {
          \draw[setag] (0.66*\x+0.66, -0.3) to[bend right=40] (0.66*\x+1.32, -0.3);
        }
        \node[leg] at (1.6,-0.95) {cada um anda uma casa para a direita};
      \end{tikzpicture}
      \end{center}
    \end{column}
    \begin{column}{0.46\textwidth}
      \begin{block}{Pense antes de passar o slide}
        \small Para pôr 40 na posição 0, os $k$ elementos que já estão lá
        precisam se deslocar.

        \medskip
        Inserindo $n$ valores, sempre no começo, quantos deslocamentos são
        feitos no total? E se $n$ dobrar, o tempo dobra?
      \end{block}
    \end{column}
  \end{columns}
\end{frame}

\begin{frame}{A resposta: medindo com \texttt{vetor\_vs\_lista.c}}
  \begin{columns}[c]
    \begin{column}{0.56\textwidth}
      \begin{center}
      \small
      \begin{tabular}{@{}rrrr@{}}
        \toprule
        $n$ & deslocamentos & vetor & lista\\
        \midrule
        10\,000  & $5{,}0 \times 10^{7}$  & 0,008\,s & 0,000\,s\\
        100\,000 & $5{,}0 \times 10^{9}$  & 0,365\,s & 0,001\,s\\
        200\,000 & $2{,}0 \times 10^{10}$ & 1,838\,s & 0,002\,s\\
        400\,000 & $8{,}0 \times 10^{10}$ & 6,109\,s & 0,003\,s\\
        \bottomrule
      \end{tabular}

      \medskip
      \scriptsize compilado com \texttt{-O2}; deslocamento com \texttt{memmove}
      \end{center}
    \end{column}
    \begin{column}{0.4\textwidth}
      \small
      $0 + 1 + \cdots + (n-1) = \frac{n(n-1)}{2}$ deslocamentos.

      \medskip
      Dobrou $n$, o vetor levou \textbf{4 vezes} mais tempo: o custo
      cresce com $n^2$.

      \medskip
      A lista não desloca ninguém: cada inserção cria um nó e acerta
      dois ponteiros.
    \end{column}
  \end{columns}
  \fonte{Cormen et al.\ (2009), §10.2~\cite{cormen2009introduction};
    medido com \texttt{vetor\_vs\_lista.c}}
\end{frame}

\begin{frame}{E o \texttt{realloc}? O bloco pode mudar de endereço}
  \begin{columns}[c]
    \begin{column}{0.5\textwidth}
      \begin{center}
      \begin{tikzpicture}
        \foreach \x/\v in {0/7,1/3,2/9,3/4} {
          \node[cel, draw=UFSGray, fill=black!5] at (0.66*\x, 1.2) {\v};
        }
        \node[rot, anchor=east] at (-0.45,1.2) {antigo};
        \foreach \x/\v in {0/7,1/3,2/9,3/4,4/,5/,6/,7/} {
          \node[cel] at (0.66*\x, 0) {\v};
        }
        \node[rot, anchor=east] at (-0.45,0) {novo};
        \draw[setag] (1,0.9) -- (1,0.3);
        \node[var] (q) at (3.6,1.9) {q};
        \draw[setar] (q) -| (1.32,1.48);
        \node[leg, text=UFSRed] at (3.6,1.35) {aponta para};
        \node[leg, text=UFSRed] at (3.6,1.05) {área liberada};
      \end{tikzpicture}
      \end{center}
    \end{column}
    \begin{column}{0.46\textwidth}
      \small
      Quando não há espaço logo depois do bloco, o \texttt{realloc}
      copia tudo para outro lugar e libera o antigo. Inserindo 200 mil
      valores, o bloco mudou de lugar \textbf{12 vezes}.

      \medskip
      Todo ponteiro guardado para dentro do vetor antigo (como
      \texttt{q = \&v[2]}) passa a apontar para memória liberada.
    \end{column}
  \end{columns}
  \fonte{ISO/IEC 9899:2018, §7.22.3.5~\cite{iso2018c17};
    medido com \texttt{vetor\_vs\_lista.c}}
\end{frame}

\begin{frame}{Definição: lista encadeada}
  \begin{block}{Lista encadeada}
    \small Uma sequência de \textbf{nós}, cada um alocado separadamente.
    Cada nó guarda um valor e o \textbf{endereço do próximo}; o último guarda
    \texttt{NULL}. Basta conhecer o endereço do primeiro nó --- o
    \textbf{início} --- para chegar a todos.
  \end{block}
  \begin{center}
  \begin{tikzpicture}[node distance=0.9cm]
    \node[var] (ini) {ini};
    \node[no, right=0.8cm of ini] (a) {40 \nodepart{two} \phantom{xx}};
    \node[no, right=of a] (b) {30 \nodepart{two} \phantom{xx}};
    \node[no, right=of b] (c) {20 \nodepart{two} \phantom{xx}};
    \node[no, right=of c] (d) {10 \nodepart{two} \phantom{xx}};
    \node[nulo, right=0.7cm of d] (n) {NULL};
    \draw[seta] (ini) -- (a);
    \draw[seta] \ptr{a} -- (b);
    \draw[seta] \ptr{b} -- (c);
    \draw[seta] \ptr{c} -- (d);
    \draw[seta] \ptr{d} -- (n);
    \node[leg, anchor=north] at ($(a.south west)!0.5!(a.one split south)$) {valor};
    \node[leg, anchor=north] at ($(a.one split south)!0.5!(a.south east)$) {prox};
  \end{tikzpicture}
  \end{center}
  \begin{exampleblock}{Intuição}
    \small Uma caça ao tesouro: cada pista diz onde está a próxima. A ordem
    está nas \emph{setas}, não na posição de cada nó na memória.
  \end{exampleblock}
  \fonte{Knuth (1997), §2.2.3~\cite{knuth1997art};
    Celes et al.\ (2004), cap.~10~\cite{celes2004introducao}}
\end{frame}

\begin{frame}{De onde veio a ideia}
  \begin{columns}[c]
    \begin{column}{0.56\textwidth}
      \small
      As listas encadeadas surgiram em 1955--56, com Allen Newell, Cliff
      Shaw e Herbert Simon, na linguagem IPL, criada para o
      \emph{Logic Theory Machine}: um programa que provava teoremas e
      precisava de estruturas que crescessem e mudassem de forma durante a
      execução.

      \medskip
      Antes disso, a memória de um programa era organizada em posições
      fixas --- como um vetor.
    \end{column}
    \begin{column}{0.4\textwidth}
      \begin{block}{O que ficou}
        \small A mesma ideia está por trás de pilhas, filas, árvores e
        grafos --- os próximos temas do curso.
      \end{block}
    \end{column}
  \end{columns}
  \fonte{Newell \& Shaw (1957)~\cite{newell1957programming};
    Knuth (1997), §2.2.3, notas históricas~\cite{knuth1997art}}
\end{frame}

\begin{frame}{Recapitulando: por que não basta um vetor}
  \begin{block}{O que aprendemos}
    \begin{itemize}
      \item num vetor, os elementos ficam lado a lado: inserir no começo
        desloca todos;
      \item inserir $n$ valores no começo custa $\approx n^2/2$
        deslocamentos;
      \item o \texttt{realloc} pode mover o bloco e invalidar ponteiros
        antigos;
      \item numa lista encadeada, cada nó mora num bloco próprio e a ordem
        é dada pelos ponteiros.
    \end{itemize}
  \end{block}
\end{frame}

% ==================================================================
\section{O nó}
% ==================================================================

\begin{frame}[fragile]{Um registro que aponta para outro do mesmo tipo}
  \begin{columns}[c]
    \begin{column}{0.46\textwidth}
\begin{lstlisting}[style=tinyC]
typedef struct no {
    int valor;
    struct no *prox;
} No;
\end{lstlisting}
    \end{column}
    \begin{column}{0.5\textwidth}
      \begin{block}{Pense antes de passar o slide}
        \small
        \begin{enumerate}
          \item Por que o campo é \texttt{struct no *prox}, e não
            \texttt{struct no prox} (sem o \texttt{*})?
          \item Por que não escrever \texttt{No *prox} lá dentro?
          \item Por que o registro tem nome (\texttt{no}) se já existe o
            apelido \texttt{No}?
        \end{enumerate}
      \end{block}
    \end{column}
  \end{columns}
\end{frame}

\begin{frame}{Definição: estrutura auto-referenciada}
  \begin{block}{Estrutura auto-referenciada}
    \small Um registro com um campo que é \textbf{ponteiro} para o próprio
    tipo do registro.
  \end{block}
  \begin{columns}[T]
    \begin{column}{0.48\textwidth}
      \small
      \textbf{Sem o \texttt{*}:} um \texttt{no} conteria um \texttt{no}, que
      conteria outro \texttt{no}\ldots{} O tamanho seria infinito, e o
      compilador recusa.

      \medskip
      \textbf{Com o \texttt{*}:} o campo guarda só um endereço --- 8 bytes,
      qualquer que seja o tamanho do nó.
    \end{column}
    \begin{column}{0.48\textwidth}
      \small
      \textbf{O nome \texttt{struct no}:} dentro das chaves, o apelido
      \texttt{No} ainda não existe --- o \texttt{typedef} só termina no
      \texttt{;}. Já \texttt{struct no} existe desde a primeira linha, como
      um tipo \emph{incompleto}: não dá para criar um, mas dá para apontar
      para um.
    \end{column}
  \end{columns}
  \fonte{Kernighan \& Ritchie (1989), §6.5 e §6.7~\cite{kernighan1989c};
    ISO/IEC 9899:2018, §6.7.2.3~\cite{iso2018c17}}
\end{frame}

\begin{frame}{Por que 16 bytes, e não 12?}
  \begin{columns}[c]
    \begin{column}{0.52\textwidth}
      \begin{center}
      \begin{tikzpicture}[x=0.36cm]
        \foreach \i in {0,...,3} \node[cel, minimum width=0.36cm] at (\i,0) {};
        \foreach \i in {4,...,7} \node[cel, minimum width=0.36cm,
          draw=UFSGray, pattern=north east lines, pattern color=UFSGray!60] at (\i,0) {};
        \foreach \i in {8,...,15} \node[cel, minimum width=0.36cm,
          draw=UFSRed, fill=UFSRed!10] at (\i,0) {};
        \draw[decorate, decoration={brace, amplitude=4pt}] (-0.5,0.35) -- (3.5,0.35)
          node[midway, above=4pt, leg] {\texttt{valor}};
        \draw[decorate, decoration={brace, amplitude=4pt}] (3.5,0.35) -- (7.5,0.35)
          node[midway, above=4pt, leg] {enchimento};
        \draw[decorate, decoration={brace, amplitude=4pt}] (7.5,0.35) -- (15.5,0.35)
          node[midway, above=4pt, leg] {\texttt{prox}};
        \foreach \i in {0,4,8,15} \node[rot, below] at (\i,-0.3) {\i};
      \end{tikzpicture}
      \end{center}
    \end{column}
    \begin{column}{0.44\textwidth}
      \small
      É o alinhamento da aula de registros: um ponteiro de 8 bytes precisa
      começar num endereço múltiplo de 8. O compilador põe 4 bytes de
      enchimento depois de \texttt{valor}.

      \medskip
      Numa lista de \texttt{int}, cada valor de 4 bytes custa
      \textbf{16}: a lista troca memória por flexibilidade.
    \end{column}
  \end{columns}
  \fonte{Bryant \& O'Hallaron (2015), §3.9.3~\cite{bryant2015computer};
    ISO/IEC 9899:2018, §6.7.2.1~\cite{iso2018c17}}
\end{frame}

\begin{frame}{Recapitulando: o nó}
  \begin{block}{O que aprendemos}
    \begin{itemize}
      \item o nó é um registro com o valor e um ponteiro para o próximo nó;
      \item o campo precisa ser ponteiro: um nó não pode conter um nó;
      \item dentro das chaves, use \texttt{struct no}: o apelido ainda não
        existe;
      \item cada nó carrega o custo do ponteiro (e do enchimento).
    \end{itemize}
  \end{block}
\end{frame}

% ==================================================================
\section{Montar e percorrer}
% ==================================================================

\begin{frame}[fragile]{Três nós montados à mão}
  \begin{columns}[c]
    \begin{column}{0.46\textwidth}
\begin{lstlisting}[style=tinyC]
No *a = malloc(sizeof *a);
No *b = malloc(sizeof *b);
No *c = malloc(sizeof *c);
a->valor = 7;  a->prox = b;
b->valor = 3;  b->prox = c;
c->valor = 9;  c->prox = NULL;
No *ini = a;
\end{lstlisting}
      \small
      \texttt{p->campo} é o mesmo que \texttt{(*p).campo}: ``siga o ponteiro
      e pegue o campo''.
    \end{column}
    \begin{column}{0.5\textwidth}
      \begin{center}
      \begin{tikzpicture}
        \node[var] (ini) at (0,1.6) {ini};
        \node[var] (va) at (1.2,1.6) {a};
        \node[var] (vb) at (2.4,1.6) {b};
        \node[var] (vc) at (3.6,1.6) {c};
        \node[leg] at (4.6,1.6) {pilha};
        \node[no] (na) at (0.9,0) {7 \nodepart{two} \phantom{xx}};
        \node[no] (nc) at (2.0,-1.1) {9 \nodepart{two} \phantom{xx}};
        \node[no] (nb) at (3.5,0.2) {3 \nodepart{two} \phantom{xx}};
        \node[nulo] (nul) at (3.7,-1.1) {NULL};
        \node[leg] at (4.6,-0.5) {heap};
        \draw[setag] (ini) -- (na);
        \draw[setag] (va) -- (na);
        \draw[setag] (vb) -- (nb);
        \draw[setag] (vc) -- (nc);
        \draw[seta] \ptr{na} -- (nb.west);
        \draw[seta] \ptr{nb} to[bend left=30] (nc.north);
        \draw[seta] \ptr{nc} -- (nul);
      \end{tikzpicture}
      \end{center}
      \scriptsize Os nós ficam onde o \texttt{malloc} quiser; a ordem
      $7, 3, 9$ está nas setas.
    \end{column}
  \end{columns}
  \fonte{King (2008), §17.5~\cite{king2008c}}
\end{frame}

\begin{frame}[fragile]{Percorrer: o que acontece com \texttt{p}?}
  \begin{columns}[c]
    \begin{column}{0.5\textwidth}
\begin{lstlisting}[style=tinyC]
for (const No *p = ini; p != NULL;
     p = p->prox)
    printf("%d -> ", p->valor);
printf("NULL\n");
\end{lstlisting}
    \end{column}
    \begin{column}{0.46\textwidth}
      \begin{block}{Pense antes de passar o slide}
        \small Com a lista $7 \to 3 \to 9$:
        \begin{enumerate}
          \item quantas vezes a condição \texttt{p != NULL} é testada?
          \item que valores \texttt{p} assume?
          \item o que acontece se a lista for vazia (\texttt{ini == NULL})?
          \item e se o último nó tiver lixo no lugar de \texttt{NULL}?
        \end{enumerate}
      \end{block}
    \end{column}
  \end{columns}
\end{frame}

\begin{frame}{A resposta: \texttt{p} anda de seta em seta}
  \begin{columns}[c]
    \begin{column}{0.48\textwidth}
      \begin{center}
      \small
      \begin{tabular}{@{}clc@{}}
        \toprule
        teste & \texttt{p} aponta para & imprime\\
        \midrule
        1 & nó 7 & \texttt{7 -> }\\
        2 & nó 3 & \texttt{3 -> }\\
        3 & nó 9 & \texttt{9 -> }\\
        4 & \texttt{NULL} & sai do laço\\
        \bottomrule
      \end{tabular}
      \end{center}
    \end{column}
    \begin{column}{0.48\textwidth}
      \small
      $n$ nós, $n+1$ testes. Com a lista vazia, o primeiro teste já falha
      e só sai \texttt{NULL}: o caso vazio funciona \emph{sem} \texttt{if}
      especial.

      \medskip
      Sem o \texttt{NULL} no fim, \texttt{p} segue um endereço de lixo:
      o programa imprime bobagem ou cai. O \texttt{NULL} é a única forma de
      saber que a lista acabou.
    \end{column}
  \end{columns}
  \fonte{Feofiloff (2009), cap.~4~\cite{feofiloff2009algoritmos};
    Kernighan \& Ritchie (1989), §6.5~\cite{kernighan1989c}}
\end{frame}

\begin{frame}[fragile]{Uma lista é vazia, ou um nó seguido de uma lista}
  \begin{columns}[c]
    \begin{column}{0.46\textwidth}
\begin{lstlisting}[style=tinyC]
int conta(const No *ini)
{
    if (ini == NULL)      /* vazia */
        return 0;
    return 1 + conta(ini->prox);
}
\end{lstlisting}
      \small A definição da lista já é recursiva --- e
      \texttt{conta} só a segue: caso base na lista vazia, chamada no
      resto.
    \end{column}
    \begin{column}{0.5\textwidth}
      \begin{center}
      \begin{tikzpicture}[node distance=0.55cm]
        \node[no] (a) {7 \nodepart{two} \phantom{xx}};
        \node[no, right=of a] (b) {3 \nodepart{two} \phantom{xx}};
        \node[no, right=of b] (c) {9 \nodepart{two} \phantom{xx}};
        \node[nulo, right=0.45cm of c] (n) {NULL};
        \draw[seta] \ptr{a} -- (b);
        \draw[seta] \ptr{b} -- (c);
        \draw[seta] \ptr{c} -- (n);
        \node[draw=UFSRed, dashed, rounded corners, fit=(b)(c)(n),
          inner sep=4pt] (r) {};
        \node[leg, text=UFSRed, below=0.05cm of r] {o resto: outra lista};
        \node[leg, text=UFSBlue, below=0.3cm of a] {um nó};
      \end{tikzpicture}
      \end{center}
      \begin{alertblock}{Lembre da aula passada}
        \small Um nó $=$ um registro na pilha. Com 300 mil nós, sem
        \texttt{-O2}, \texttt{conta} estoura a pilha.
      \end{alertblock}
    \end{column}
  \end{columns}
  \fonte{Feofiloff (2009), cap.~4~\cite{feofiloff2009algoritmos};
    medido com \texttt{lista.c}}
\end{frame}

\begin{frame}{Recapitulando: montar e percorrer}
  \begin{block}{O que aprendemos}
    \begin{itemize}
      \item \texttt{p->campo} segue o ponteiro e acessa o campo;
      \item percorrer é andar \texttt{p = p->prox} até \texttt{p == NULL};
      \item o mesmo laço trata a lista vazia, sem caso especial;
      \item a lista é uma definição recursiva --- mas a recursão gasta um
        registro por nó.
    \end{itemize}
  \end{block}
\end{frame}

% ==================================================================
\section{Inserir no início}
% ==================================================================

\begin{frame}{Por que a ordem saiu invertida?}
  \begin{columns}[c]
    \begin{column}{0.5\textwidth}
      \small
      \texttt{main} chamou, nesta ordem:\\[0.3em]
      \texttt{\ \ insere\_inicio(lista, 10)}\\
      \texttt{\ \ insere\_inicio(lista, 20)}\\
      \texttt{\ \ insere\_inicio(lista, 30)}\\
      \texttt{\ \ insere\_inicio(lista, 40)}

      \medskip
      e \texttt{imprime} mostrou \texttt{40 -> 30 -> 20 -> 10}.
    \end{column}
    \begin{column}{0.46\textwidth}
      \begin{block}{Pense antes de passar o slide}
        \small
        \begin{enumerate}
          \item Desenhe a lista depois de cada chamada.
          \item Por que \texttt{insere\_inicio} \emph{devolve} um ponteiro?
          \item O que acontece se o \texttt{main} escrever só
            \texttt{insere\_inicio(lista, 50);}, sem atribuir?
        \end{enumerate}
      \end{block}
    \end{column}
  \end{columns}
\end{frame}

\begin{frame}{Inserir no início: dois ponteiros, nesta ordem}
  \begin{columns}[c]
    \begin{column}{0.56\textwidth}
  \begin{tikzpicture}[node distance=0.55cm]
    \node[leg, anchor=west] at (-0.4,0.75) {\textbf{1.} \texttt{novo->prox = ini;}};
    \node[var] (ini1) at (0,0) {ini};
    \node[no, right=0.55cm of ini1] (a1) {20 \nodepart{two} \phantom{xx}};
    \node[no, right=of a1] (b1) {10 \nodepart{two} \phantom{xx}};
    \node[nulo, right=0.45cm of b1] (n1) {NULL};
    \node[var] (nv1) at (0,-0.85) {novo};
    \node[nonovo, right=0.55cm of nv1] (x1) {30 \nodepart{two} \phantom{xx}};
    \draw[seta] (ini1) -- (a1);
    \draw[seta] \ptr{a1} -- (b1);
    \draw[seta] \ptr{b1} -- (n1);
    \draw[setag] (nv1) -- (x1);
    \draw[setar] \ptr{x1} -- ++(0.5,0) |- ($(a1.south east)!0.5!(a1.south west)+(0.25,-0.18)$) -- ++(0,0.18);
    \node[leg, anchor=west] at (-0.4,-1.75) {\textbf{2.} \texttt{return novo;} \quad e \texttt{lista = \ldots} no \texttt{main}};
    \node[var] (ini2) at (0,-2.5) {lista};
    \node[nonovo, right=0.55cm of ini2] (x2) {30 \nodepart{two} \phantom{xx}};
    \node[no, right=of x2] (a2) {20 \nodepart{two} \phantom{xx}};
    \node[no, right=of a2] (b2) {10 \nodepart{two} \phantom{xx}};
    \node[nulo, right=0.45cm of b2] (n2) {NULL};
    \draw[setar] (ini2) -- (x2);
    \draw[seta] \ptr{x2} -- (a2);
    \draw[seta] \ptr{a2} -- (b2);
    \draw[seta] \ptr{b2} -- (n2);
  \end{tikzpicture}
    \end{column}
    \begin{column}{0.4\textwidth}
      \small
      Na ordem contrária --- primeiro mudar o início, depois ligar --- o
      endereço do antigo primeiro nó se perderia, e com ele a lista
      inteira.

      \medskip
      São sempre os mesmos passos, com 3 ou 3 milhões de nós: custo
      \textbf{constante}.
    \end{column}
  \end{columns}
  \fonte{Celes et al.\ (2004), cap.~10~\cite{celes2004introducao};
    Cormen et al.\ (2009), §10.2~\cite{cormen2009introduction}}
\end{frame}

\begin{frame}[fragile]{A resposta: \texttt{ini} é uma cópia}
  \begin{columns}[T]
    \begin{column}{0.48\textwidth}
\begin{lstlisting}[style=tinyC]
No *lista = NULL;
insere_inicio(lista, 50);  /* ERRADO */
imprime(lista);            /* NULL */
\end{lstlisting}
      \small A função recebe uma \textbf{cópia} de \texttt{lista}. Mudar
      \texttt{ini} lá dentro não muda \texttt{lista} aqui fora --- a mesma
      passagem por valor da aula de ponteiros.
    \end{column}
    \begin{column}{0.48\textwidth}
      \begin{block}{Duas consequências}
        \small
        \begin{itemize}
          \item \texttt{lista} continua \texttt{NULL};
          \item o nó criado fica no heap sem ninguém apontando para ele:
            um \textbf{vazamento}.
        \end{itemize}
      \end{block}
      \small Por isso a função devolve o novo início, e quem chama
      \emph{atribui}: \texttt{lista = insere\_inicio(lista, 50);}
    \end{column}
  \end{columns}
  \fonte{Kernighan \& Ritchie (1989), §5.2~\cite{kernighan1989c};
    King (2008), §17.5~\cite{king2008c}}
\end{frame}

\begin{frame}{Recapitulando: inserir no início}
  \begin{block}{O que aprendemos}
    \begin{itemize}
      \item o novo nó aponta para o antigo início; depois ele vira o início;
      \item a ordem dos dois passos importa: ligar antes de mudar o início;
      \item quem chama precisa guardar o ponteiro devolvido;
      \item o custo não depende do tamanho da lista --- e por isso a ordem
        sai invertida.
    \end{itemize}
  \end{block}
\end{frame}

% ==================================================================
\section{Liberar a lista}
% ==================================================================

\begin{frame}[fragile]{Um \texttt{libera} mais curto}
  \begin{columns}[c]
    \begin{column}{0.52\textwidth}
\begin{lstlisting}[style=tinyC]
/* libera_errada.c */
void libera(No *ini)
{
    for (No *p = ini; p != NULL;
         p = p->prox)
        free(p);
}
\end{lstlisting}
      \small Três linhas, o mesmo laço do \texttt{imprime}, com
      \texttt{free} no lugar do \texttt{printf}.
    \end{column}
    \begin{column}{0.44\textwidth}
      \begin{block}{Pense antes de passar o slide}
        \small
        \begin{enumerate}
          \item Esta versão libera os quatro nós?
          \item Em que momento é lido \texttt{p->prox}: antes ou depois do
            \texttt{free(p)}?
          \item O programa vai funcionar, cair, ou depende?
        \end{enumerate}
      \end{block}
    \end{column}
  \end{columns}
\end{frame}

\begin{frame}[fragile]{Execute e compare}
\begin{lstlisting}[style=shellStyle]
$ gcc -Wall -Wextra libera_errada.c -o libera_errada     # nenhum aviso
$ ./libera_errada
Abort trap: 6
$ gcc -g -fsanitize=address libera_errada.c -o libera_errada-asan
$ ./libera_errada-asan
==4117==ERROR: AddressSanitizer: heap-use-after-free on address 0x602000000158
READ of size 8 at 0x602000000158 thread T0
    #0 0x000102040950 in libera libera_errada.c:22
0x602000000158 is located 8 bytes inside of 16-byte region
freed by thread T0 here:
    #1 0x000102040920 in libera libera_errada.c:23
\end{lstlisting}
  \begin{center}
    \scriptsize\textcolor{UFSGray}{No macOS, o alocador derrubou o programa.
    Em outro sistema, o mesmo código pode cair de outra forma --- ou parecer
    funcionar.}
  \end{center}
\end{frame}

\begin{frame}{A resposta: o \texttt{p->prox} vem depois do \texttt{free}}
  \begin{columns}[c]
    \begin{column}{0.48\textwidth}
      \begin{center}
      \begin{tikzpicture}[node distance=0.6cm]
        \node[nolivre] (a) {\phantom{0}\nodepart{two}\phantom{xx}};
        \node[no, right=of a] (b) {30 \nodepart{two} \phantom{xx}};
        \node[no, right=of b] (c) {20 \nodepart{two} \phantom{xx}};
        \node[var, above=0.5cm of a] (p) {p};
        \draw[setag] (p) -- (a);
        \draw[setar, dashed] \ptr{a} -- node[above, leg, text=UFSRed] {?} (b);
        \draw[seta] \ptr{b} -- (c);
        \node[leg, below=0.2cm of a] {já liberado};
      \end{tikzpicture}
      \end{center}
      \small Depois de \texttt{free(p)}, o bloco volta para o alocador,
      que pode reescrevê-lo. ``8 bytes dentro de uma região de 16'' é
      exatamente o campo \texttt{prox}.
    \end{column}
    \begin{column}{0.48\textwidth}
      \begin{block}{A correção: guarde o caminho antes}
        \small
        \texttt{No *prox = ini->prox;}\\
        \texttt{free(ini);}\\
        \texttt{ini = prox;}
      \end{block}
      \small Usar um bloco depois do \texttt{free} é comportamento
      indefinido: o padrão não promete nada. O \emph{sanitizer} transforma
      ``depende'' em uma mensagem com a linha do erro.
    \end{column}
  \end{columns}
  \fonte{ISO/IEC 9899:2018, §7.22.3.3~\cite{iso2018c17};
    Serebryany et al.\ (2012)~\cite{serebryany2012asan}}
\end{frame}

\begin{frame}{Recapitulando: liberar a lista}
  \begin{block}{O que aprendemos}
    \begin{itemize}
      \item cada \texttt{malloc} de nó precisa do seu \texttt{free}: um laço
        que visita todos;
      \item depois do \texttt{free}, o nó não pode mais ser lido --- nem o
        \texttt{prox};
      \item guarde o próximo numa variável \emph{antes} de liberar o atual;
      \item o compilador não avisa; use \texttt{-fsanitize=address} para
        encontrar o erro.
    \end{itemize}
  \end{block}
\end{frame}

% ==================================================================
\section{Vetor ou lista?}
% ==================================================================

\begin{frame}{Cada estrutura é boa em uma coisa}
  \begin{center}
    \small
    \begin{tabular}{@{}lcc@{}}
      \toprule
      operação & vetor & lista encadeada\\
      \midrule
      acessar o $k$-ésimo elemento & imediato (\texttt{v[k]}) & percorre $k$ nós\\
      inserir no começo & desloca $n$ & \textbf{constante}\\
      inserir no fim & constante\textsuperscript{*} & percorre $n$ nós\\
      memória por \texttt{int} & 4 bytes & 16 bytes\\
      elementos vizinhos na memória & sim & não\\
      \bottomrule
    \end{tabular}

    \smallskip
    \scriptsize\textsuperscript{*} em média, com \texttt{realloc} dobrando a
    capacidade
  \end{center}
  \small
  Elementos vizinhos aproveitam melhor a \emph{cache} do processador: somar
  10 milhões de \texttt{int} levou 0,005\,s no vetor e 0,025\,s na lista ---
  5 vezes mais, mesmo com os nós alocados em sequência.
  \fonte{Cormen et al.\ (2009), §10.2~\cite{cormen2009introduction};
    Bryant \& O'Hallaron (2015), §6.2~\cite{bryant2015computer};
    medido com \texttt{percorre.c}}
\end{frame}

\begin{frame}{Recapitulando: vetor ou lista?}
  \begin{block}{O que aprendemos}
    \begin{itemize}
      \item vetor: acesso direto por índice, mas inserir fora do fim
        desloca elementos;
      \item lista: inserir no começo é barato, mas chegar ao $k$-ésimo
        exige percorrer;
      \item a lista gasta mais memória por elemento e usa pior a cache;
      \item a escolha depende das operações que o programa mais faz.
    \end{itemize}
  \end{block}
\end{frame}

% ==================================================================
\section{Mãos à obra}
% ==================================================================

\begin{frame}{Altere o programa}
  \small
  \begin{block}{Tarefa 1}
    Em \texttt{imprime}, mostre também o endereço de cada nó
    (\texttt{\%p}, com \texttt{(void *) p}). Os nós estão lado a lado? A
    distância entre eles é de 16 bytes?
  \end{block}
  \begin{block}{Tarefa 2}
    Troque o \texttt{4} do laço do \texttt{main} por \texttt{1000000} e tire o
    \texttt{imprime}. Preveja: \texttt{conta} funciona sem \texttt{-O2}? E
    com? Depois reescreva \texttt{conta} com um laço.
  \end{block}
  \begin{block}{Tarefa 3}
    Escreva \texttt{int soma(const No *ini)} duas vezes: com laço e
    recursiva. Teste com a lista vazia.
  \end{block}
\end{frame}

\begin{frame}{Exercícios (1/2)}
  \footnotesize
  \begin{block}{Exercício 1 --- o maior valor}
    \texttt{int maximo(const No *ini)} devolve o maior valor da lista.
    O que fazer quando a lista é vazia? Escolha e documente no comentário.
  \end{block}
  \begin{block}{Exercício 2 --- do vetor para a lista}
    \texttt{No *de\_vetor(const int *v, int n)} cria uma lista com os valores
    de \texttt{v} \textbf{na mesma ordem}, usando só \texttt{insere\_inicio}.
    Dica: por qual ponta do vetor começar?
  \end{block}
  \begin{block}{Exercício 3 --- imprimir de trás para frente}
    \texttt{void imprime\_invertida(const No *ini)}, recursiva, sem criar
    outra lista. Onde fica o \texttt{printf}: na ida ou na volta?
  \end{block}
  \fonte{Feofiloff (2009), cap.~4~\cite{feofiloff2009algoritmos};
    Celes et al.\ (2004), cap.~10~\cite{celes2004introducao}}
\end{frame}

\begin{frame}{Exercícios (2/2)}
  \footnotesize
  \begin{block}{Exercício 4 --- listas iguais}
    \texttt{int iguais(const No *a, const No *b)} devolve 1 se as duas listas
    têm os mesmos valores, na mesma ordem, e 0 caso contrário. Cuidado com
    listas de tamanhos diferentes.
  \end{block}
  \begin{block}{Exercício 5 --- uma cópia independente}
    \texttt{No *copia(const No *ini)} devolve uma lista nova, na mesma ordem,
    com nós próprios. Libere a original e mostre que a cópia continua
    funcionando.
  \end{block}
  \begin{block}{Exercício 6 --- \texttt{k}-ésimo}
    \texttt{const No *kesimo(const No *ini, int k)} devolve o nó da posição
    \texttt{k} (a partir de 0), ou \texttt{NULL} se ela não existir.
    Quantos nós são visitados para \texttt{k = n - 1}?
  \end{block}
  \fonte{King (2008), §17.5~\cite{king2008c};
    Tenenbaum et al.\ (1995), cap.~4~\cite{tenenbaum1995estruturas}}
\end{frame}

\begin{frame}{Desafio}
  \begin{block}{Contatos sem limite, com nomes do tamanho certo}
    \small Leia nomes da entrada padrão, um por linha, até o fim do arquivo,
    e guarde-os numa lista encadeada de
    \begin{center}
      \texttt{typedef struct contato \{ char *nome; struct contato *prox; \} Contato;}
    \end{center}
    \vspace{-0.6em}
    \begin{itemize}
      \item cada nome é lido num vetor local de 100 posições e
        \textbf{copiado} para um bloco do \texttt{malloc} com o tamanho exato;
      \item imprima os nomes em ordem inversa à da leitura e, depois, na
        ordem da leitura;
      \item no fim, libere tudo: dois \texttt{free} por nó --- em que ordem?
    \end{itemize}
    Com \texttt{-fsanitize=address}, nenhum erro e nenhum vazamento.
  \end{block}
  \fonte{King (2008), §17.2 e §17.5~\cite{king2008c}}
\end{frame}

\begin{frame}{Fechamento}
  \begin{columns}[T]
    \begin{column}{0.5\textwidth}
      \textbf{O que mudou hoje}\\[0.3em]
      \small A sequência deixou de depender da posição na memória: a ordem
      agora está nos ponteiros. Inserir no começo ficou barato --- com 3 ou
      3 milhões de nós.
    \end{column}
    \begin{column}{0.5\textwidth}
      \textbf{O preço}\\[0.3em]
      \small Mais memória por elemento, nada de acesso por índice, e um
      \texttt{malloc} e um \texttt{free} para cada nó.
    \end{column}
  \end{columns}
  \begin{block}{Para levar}
    \begin{itemize}
      \small
      \item o nó é um registro com valor e \texttt{struct no *prox};
      \item percorrer: \texttt{p = p->prox} até \texttt{NULL};
      \item inserir no início: ligar o novo nó, depois devolvê-lo como início;
      \item liberar: guardar o próximo antes do \texttt{free}.
    \end{itemize}
  \end{block}
  \small\textcolor{UFSGray}{Próxima aula: busca, remoção e inserção em
  qualquer posição.}
\end{frame}

\section*{Referências}
\begin{frame}[allowframebreaks]{Referências}
  \nocite{*}
  \bibliographystyle{plain}
  \bibliography{referencias}
\end{frame}

\end{document}
